This section describes the mid-training procedure. We do not pre-train our own model from scratch, but start from an existing pretrained base model. We further mid-train it on the synthetic data introduced in Section~\ref{sec:training-data}.

\paragraph{Base Model}
We compared three open-weight families of small language models as starting points: Qwen3~\citep{qwen3}, Gemma3~\citep{gemma3}, and SmolLM3~\citep{smollm3}. Selection was based on early runs evaluated on a held-out QA slice. Qwen3 produced the strongest result at fixed compute and was retained for both released sizes. Our final checkpoints are obtained by mid-training \texttt{Qwen3-0.6B-Base} and \texttt{Qwen3-1.7B-Base}.

\paragraph{Training Objective and Format}
We train via supervised fine-tuning, with the loss applied only to the response tokens. The complete prompt/response format is shown in Figure~\ref{fig:trace-training-example} (Appendix~\ref{sec:prompt-trace}). The prompt is a question together with the context passages in random order, each tagged with a numeric source identifier; this format is identical to the format used at evaluation time, so no train-test mismatch is introduced. The response is a reasoning trace, generated and formatted as described in Section~\ref{sec:synth-traces}; note that the final answer and the \textsc{Answerable}/\textsc{Unanswerable} verdict are included in the reasoning trace. To make boundaries between prompt elements (query, context passages) and response sections (query analysis, source analysis, etc.) explicit, we delimit them by a small set of additional special tokens; their embeddings are added to the base model and initialised from the mean of the subword embeddings of their natural-language names~\citep{hewitt2021vocab}.

\paragraph{Data Mixing}
As described in Section~\ref{sec:training-data}, our training corpus contains three main subsets: (1) single-hop, (2) multi-hop single-context, and (3) multi-hop multi-context. The single-hop subset is roughly an order of magnitude larger than each of the multi-hop subsets, but multi-hop examples are the ones that actually exercise the reasoning capability we want the model to develop. We therefore oversample both multi-hop subsets: each multi-hop example is shown three times within a single epoch, while every single-hop example is shown once. This consistently improves multi-hop benchmark accuracy without measurable regression on single-hop. We also evaluated a curriculum schedule that trains only on single-hop examples at the beginning and starts mixing in multi-hop data from a fixed step onwards, but observed no measurable difference compared to the static mixture.

\paragraph{Implementation Details}
OCC-RAG-0.6B and OCC-RAG-1.7B were both trained on approximately $9 \times 10^9$ tokens, respectively, taking approximately $17$ and $28$ wall-clock hours each on $8$ NVIDIA H100 (80~GB) GPUs. Full training hyperparameters and the distributed-training configuration are listed in Appendix~\ref{appendix:training-hyperparameters}.
